\documentclass{beamer}
% \usepackage{color}
\usepackage{graphicx}
\usepackage{setspace}
\setstretch{1.5}
% \usepackage[all,cmtip]{xy}
\usepackage{tikz}
\usepackage{tikz-cd}
\usepackage{qtree}
\usepackage{stmaryrd}
% \usepackage{diagxy}
\usepackage{amsfonts}
\usepackage{amssymb}
\usepackage{amsmath}
\usepackage{amscd}
\usepackage{amsthm}
% \usepackage{mathpartir}
\usepackage{bbm}
%\usepackage{pictexwd,dcpic}
% \usepackage[mathcal]{euscript}
\usepackage[utf8]{inputenc}
%\usepackage[T1]{fontenc}
% \usepackage{lmodern}
\title{reality and equivalence \\ in algebraic quantum theory}
\author{Hans Halvorson}
\date{Nagoya Winter Workshop \\ March 11, 2015}
% \usetheme{Pittsburgh}
\usetheme{Singapore}
\setbeamertemplate{navigation symbols}{}%remove navigation symbols
% \usetheme{Singapore}
\addtobeamertemplate{navigation symbols}{}{%
    \usebeamerfont{footline}%
    \usebeamercolor[fg]{footline}%
    \hspace{1em}%
    \insertframenumber/\inserttotalframenumber
}
\usepackage{eucal}
\def\2{\mathcal}
\def\7{\mathfrak}


\begin{document}

\begin{frame}

\maketitle

\end{frame}


\begin{frame}{overview}

\begin{enumerate}

\item preliminaries

\item inequivalent representations

\item equivalent QFTs

\item same physics on all spacetimes


\end{enumerate}



\end{frame}


\begin{frame}{what does equivalence have to do with reality?}

  \begin{centering} \textcolor{blue}{\bf mantra:} reality is invariant
    relative to equivalence of true descriptions of the
    world. \end{centering}

% coordinates in differential geometry

% electrostatic potential

% Japanese vs. English descriptions


\end{frame}



\begin{frame}{what sorts of things can be equivalent?}

\begin{itemize}
\item sentences
\item descriptions
\item theories
\begin{itemize}
\item sentences?
\item models?
\item functors?
\end{itemize}
\item representations
\item Lagrangians

\end{itemize}


\end{frame}


\begin{frame}{inequivalent quantizations}

  \textcolor{green}{OBJECTIVE:} describe a quantum system
  corresponding to a classical system with configuration space $X$.
  

  \bigskip

  \textcolor{red}{OBSTACLE:} embarrassment of riches --- if $X$ is not
  simply connected, then there are inequivalent quantum theories
  corresponding to $X$.

\end{frame}

\begin{frame}{from quantization to representation}

\begin{itemize}
\item quantization: given configuration space $X$, find an
  corresponding quantum theory (hilbert space and operators satisfying
  certain algebraic relations).
\item Landsman's maneuver: there is a unique $C^*$-algebra $A[X]$ such
  that Hilbert space representations of $A[X]$ are in one-to-one
  correspondence with quantizations of $X$.
\end{itemize}

\end{frame}


\begin{frame}{algebraic imperialism (I.E. Segal)}

  \begin{quote} ``The physical content of the theory is completely
    contained in the abstract algebra.'' \end{quote}

  \bigskip \textbf{Observables:} self-adjoint operators in $A$.

\bigskip \textbf{States:} positive, linear, normalized maps from
observables to $\mathbb{R}$.


\bigskip \textbf{Dynamics:} one-parameter group of automorphisms of
$A$.




\end{frame}

\begin{frame}{hilbert space conservatism} 


  \textbf{Representation:} a module $\pi :A\times H\to H$

  \bigskip \textbf{Fact:} A representation $\pi :A\times H\to H$
  induces a map $\pi ^*:T_1(H)\to S(A)$, given by
$$ \pi ^*(d)(a) \: = \: \mathsf{Tr}(d\pi (a)) .$$


\bigskip \textbf{Gelfand-Naimark-Segal Theorem:} $\forall \omega \in
S(A)$, there is a Hilbert space $H_\omega$, an operator $d\in
T_1(H_\omega )$, and a representation $\pi _\omega :A\times
H_\omega\to H_\omega$ such that $\omega =\pi _\omega ^*(d)$.

\end{frame}



\begin{frame}{unitarily equivalent representations}

  A unitary equivalence $w:(H,\pi )\to (K,\phi )$ is a unitary map
  $w:H\to K$ that makes the following diagram commute:

$$
\begin{tikzcd}[ampersand replacement=\&]
  A \times H \arrow{r}{\pi } \arrow{d}[left]{1\times w } \& H \arrow{d}[right]{w} \\
  A \times K \arrow{r}[below]{\phi } \& K
\end{tikzcd}
$$

\end{frame}


\begin{frame}{disregarding inequivalent representations}

  \textbf{Fell's theorem:} for any $d\in T_1(H)$, operators $a_1,\dots
  ,a_n\in A$, and $\epsilon >0$, there is a $d'\in T_1(K)$ such that
$$ \bigl| \mathsf{Tr}(d\pi(a_i))- \mathsf{Tr}(d'\phi (a_i)) \bigr| \: < \: \epsilon .$$

\bigskip Claim (Segal): any two representations of $A[X]$ are
physically equivalent.


\end{frame}


\begin{frame}{regarding inequivalent representations}

  Claim (Ruetsche): physically equivalent representations are
  unitarily equivalent.

  \bigskip $\Rightarrow$ for non simply connected $X$, there are
  physically inequivalent representations of $A[X]$.



\end{frame}






\begin{frame}{what's the issue?}

\begin{enumerate}
\item hilbert space conservatism: there is one correct representation

  {\small $\Rightarrow$ choosing a representation involves a
    further theoretical commitment}

\item algebraic imperialism: all physical information is in the
  abstract algebra of observables

  \smallskip {\small $\Rightarrow$ choosing a representation entails no further theoretical commitment}
\end{enumerate}

\end{frame}


% \begin{frame}{algebraic quantum theory}

%   A $\mathbf{C^*}$\textbf{-algebra} is:
% \begin{itemize}
% \item An algebra over $\mathbb{C}$, i.e.\ a ring with complex scalar
%   multiplication;
% \item with a conjugate-linear, product-reversing involution $a\mapsto
%   a^*$,
% \item with a norm $\|\cdot \| :A\to \mathbb{R}^+$ that satisfies the
%   $C^*$-property $\| a^*a\| = \| a\| ^2$,
% \item and that is complete relative to this norm.
% \end{itemize}

% \bigskip \textbf{Example:} The algebra of bounded linear operators on
% a Hilbert space $H$. 

% \end{frame}



\begin{frame}{problems for hilbert space conservatism}

\begin{itemize}
\item underdetermination: what reason for choosing a representation?
\item what is gained by choosing a representation?
\item what is lost by choosing a representation?
\item what kind of epistemic commitment is it to choose a
  representation? 
\end{itemize}

\end{frame}



\begin{frame}{Baker's conundrum}

\begin{enumerate}
\item physical equivalence $\Leftrightarrow$ intertranslatability
  (Poincar\'e, Glymour, Quine, etc.)
\item fact: there are inequivalent representations that are
  intertranslatable.
\item conclusion: there are inequivalent representations that are
  physically equivalent.
\end{enumerate}


\end{frame}


% now for the free Fermi field

%% draw picture A(O) \subseteq F(O)

%% A(O) = gauge invariant elements of F(O)

%% A = inductive limit of all A(O), same with F

%% two subspaces H_{-} and H_{+} that are invariant under action of A

%% Fact: \pi _{-}(A) and \pi _{+}(A) are unitarily inequivalent, yet isomorphic

%% Puzzle: Why not just throw away one of \pi _{-} or \pi _{+} ?



% \begin{frame}{Example 1: free Majorana field}

% \begin{itemize}
% \item gauge transformation: $\gamma :F\to F$

% %% puzzle: if A is invariant under \gamma , then how does \gamma change sectors?

% %% perhaps Morsella explains the local endomorphisms?

% \item observables are fixed points in $F$ under action of $\gamma$
% \item two inequivalent representations: $\pi _+$ vacuum, $\pi _{-}$ charged
% \item localized automorphisms of $A$
% $$ \pi _{-}(a) \: \cong \: \pi _{+}(\theta (a)) .$$
% \item $H_{+}$ and $H_{-}$ are dynamical islands
% \item field operators map $H_{+}$ to $H_{-}$.
% \end{itemize}


% %% open question: what does the free Fermi field model explain?




% \end{frame}




\begin{frame}{Buridan's Ass meets QFT}

\begin{enumerate}
\item The system can be represented by $\pi$ or by $\phi$, but not by
  both.
%% must choose a state
%% does choosing a state require choosing a representation?
%% does choosing a state give a representation for free?
\item If the world is represented by $\pi$, it could equally well be
  represented by $\phi$.
%% the two representations are isomorphic in all ways!
%% here I agree with Baker.  These two reps have the same representational power
\end{enumerate}

\end{frame}

\begin{frame}{the dilemma restated}

\begin{enumerate}
%% Ruetsche and Earman
\item break the symmetry between the two representations
%% Baker
\item deny that there is any choice to be made
%% suppose that Buridan's ass accepts PII
%% he concludes that he is seeing double
%% .. or suffering from an illusion
%% the correct description is one in which there is just one haystack
\end{enumerate}

\bigskip
common theme: ``freedom to choose a representation'' is
theoretically insignificant

% compare with: ``freedom to choose a coordinate chart'' is
% theoretically insignificant

\end{frame}



\begin{frame}{equivalent nets}

  Let $\mathbb{K}$ be a family of ``nice'' subsets of Minkowski
  spacetime.  A \textbf{net of observable algebras} is a map
  $A:\mathbb{K}\to \mathbf{CStar}$ satisfying:

\begin{itemize}
\item isotony: if $O_1\subseteq O_2$ then $A(O_1)\subseteq A(O_2)$
\item microcausality: if $O_1\perp O_2$ then $[A(O_1),A(O_2)]=\{ 0\}$
\end{itemize}

\end{frame}

\begin{frame}{equivalent nets}

  Let $A,B:\mathbb{K}\to \mathbf{CStar}$ be nets of observable
  algebras over Minkowski spacetime.  A morphism $\alpha :A\to B$ is
  defined as a family $\alpha _O:A(O)\to B(O)$ of $*$-homomorphisms
  such that for any $O_1\subseteq O_2$, the following diagram
  commutes:

$$
\begin{tikzcd}[ampersand replacement=\&]
  A(O_2) \arrow{r}{\alpha _{O_2}}    \& B(O_2) \\
  A(O_1) \arrow{u}[left]{ } \arrow{r}[below]{\alpha _{O_1}} \& B(O_1) \arrow{u}[right]{ } 
\end{tikzcd}
$$

\end{frame}


\begin{frame}{equivalent nets}

  \textcolor{purple}{question:} how much information about
  \textbf{dynamics} is contained in the net?

  \bigskip example: let $A_m$ be the net for the free Bose field of
  mass $m>0$.  Are $A_m$ and $A_{m'}$ equivalent nets when $m\neq m'$?

  \bigskip (It is known that $A_m(O)\cong A_{m'}(O)$ for all $O$.
  ``No'' would indicate that the net contains information about the
  dynamics.)


\end{frame}

\begin{frame}{locally covariant QFT}


  Let $\mathbf{Loc}$ be the category whose objects are globally
  hyperbolic manifolds, with Lorentzian metric, orientation, and time
  orientation; and whose arrows are causally convex embeddings.

  \bigskip A \textcolor{purple}{locally covariant QFT} is a covariant
  functor $A$ from $\mathbf{Loc}$ to $\mathbf{CStar}$.



\end{frame}


\begin{frame}{equivalent QFTs}

  \textbf{definition:} a \textcolor{purple}{morphism} of locally
  covariant QFTs is a natural transformation $\alpha :A\to B$.  This
  morphism is called an \textcolor{purple}{equivalence} just in case
  each $\alpha _M$ is an isomorphism.

\end{frame}


\begin{frame}{equivalent QFTs}

  \textcolor{purple}{question:} suppose $A,B$ are locally covariant
  QFTs such that $A(M)\cong B(M)$, for all spacetimes $M$.  Does it
  follow that $A$ and $B$ are equivalent?

  \bigskip \textcolor{purple}{answer:} no (Fewster)




\end{frame}


\begin{frame}{same physics on all spacetimes}

  \textbf{constructive QFT:} same construction recipe, different
  background spacetimes

  \bigskip \textbf{reverse engineering:} given a locally covariant QFT
  $A$, when should we say that $A$ describe the ``same physics'' on
  all spacetimes?





\end{frame}


\begin{frame}{same physics on all spacetimes}

  \textbf{weak necessary condition:} for two QFTs $A,B$, if both
  describe the same physics on all spacetimes, and $A(M)\cong B(M)$
  for some spacetime $M$, then $A\cong B$.


\end{frame}




\end{document}